\section{Reproducibility and Data Supplement}

This appendix expands the reproducibility checklist into a data and configuration supplement. It records the files used in each stage, the cleaned WACC panel coverage, the main model settings, and the additional diagnostic outputs used to interpret the results. The purpose is to make the IEEE version self-contained while preserving the original three-stage workflow.

\subsection{Repository and Execution Environment}

\begingroup\raggedright
All experiments follow the repository workflow used in the original submission. The primary notebooks are \texttt{Price\_Prediction.ipynb}, \texttt{Prediction\_Inflation.ipynb}, and \texttt{Impact\_FossilFuel.ipynb}. The revision diagnostics are generated by \texttt{code/revision\_experiments.py}. The versioned source and corrected diagnostics are available in the \href{https://github.com/GreenComp-ERC/ImpactOfFossilFeulPrice}{organization repository}. Experiments were run on macOS with an Apple arm64 10-core CPU and 16GB RAM. The original notebooks specify Python, AutoGluon, TensorFlow/Keras, and scikit-learn. The corrected diagnostic script depends on pandas, NumPy, and openpyxl; its rerun environment and checks are recorded in the versioned repository.\par\endgroup

\begin{table}[!ht]
\centering
\caption{Files Included in the IEEE Source Folder}
\label{tab:appendix-file-inventory}
\scriptsize
\begin{tabular}{P{0.24\columnwidth}P{0.28\columnwidth}P{0.36\columnwidth}}
\toprule
\textbf{Folder} & \textbf{File} & \textbf{Purpose} \\
\midrule
\texttt{figures/} & \texttt{LSTM.png} & Stage 1 architecture figure \\
\texttt{figures/} & \texttt{Accuracy.png} & Oil price forecast comparison \\
\texttt{figures/} & \texttt{figure.png} & Stage 3 regional/project metrics \\
\texttt{figures/} & Tensor framework PDF & First-page tensor workflow figure \\
\texttt{figures/} & Methodology relations PDF & Three-stage method and model-relation figure \\
\texttt{data/} & Stage 1 metrics & Forecast metric table \\
\texttt{data/} & Stage 3 holdout & 2033-2035 holdout checks \\
\texttt{data/} & Rolling metrics & Rolling-origin checks \\
\texttt{data/} & Shock elasticities & 10\% price-shock sensitivity \\
\texttt{data/} & Granger diagnostics & Transmission diagnostics \\
\bottomrule
\end{tabular}
\end{table}

\subsection{Data Sources and Cleaning}

Stage 1 uses historical coal, gas, and oil prices from World Bank commodity data. Stage 2 uses macroeconomic indicators, including consumer price inflation, gross domestic product growth, unemployment, real interest rate, money supply, and trade-balance indicators. Stage 3 combines predicted fuel prices and inflation with WACC projections derived from Calcaterra et al. The cleaned WACC inputs keep a common monthly horizon from 2026-01 to 2035-12 for China, Europe, and the Middle East.

\begin{table}[!ht]
\centering
\caption{Cleaned WACC Panel Coverage}
\label{tab:appendix-cleaned-wacc}
\scriptsize
\begin{tabular}{lrrrr}
\toprule
\textbf{Region} & \textbf{Raw} & \textbf{Clean} & \textbf{Horizon} & \textbf{Mean Fix} \\
\midrule
China & 120 & 120 & 2026-2035 & 0.000 \\
Europe & 120 & 120 & 2026-2035 & 1.090 \\
Middle East & 180 & 120 & 2026-2035 & 1.031 \\
\bottomrule
\end{tabular}
\vspace{2pt}

\raggedright\scriptsize
\textit{Note:} ``Mean Fix'' is the mean absolute correction to the inflation column, in percentage points. Missing values after cleaning are zero for all regions. The original Europe and Middle East workbooks used China's inflation series; cleaning replaces it with the respective regional forecasts. The original AutoGluon fit summaries remain archived separately from the cleaned-panel diagnostics.
\end{table}

\begin{table}[!ht]
\centering
\caption{Core Variables Used by Stage}
\label{tab:appendix-core-variables}
\scriptsize
\begin{tabular}{P{0.23\columnwidth}P{0.35\columnwidth}P{0.30\columnwidth}}
\toprule
\textbf{Stage} & \textbf{Variables} & \textbf{Role} \\
\midrule
Stage 1 & Fuel price, month sine/cosine, rolling averages & Price forecasting input and target \\
Stage 2 & CPI inflation, GDP growth, unemployment, real interest rate, money supply, trade indicators & Inflation forecasting input and target \\
Stage 3 & Coal, gas, oil price; inflation; WACC by project type & WACC sensitivity estimation \\
Diagnostics & Log price differences; differenced inflation and interest rates & Transmission mechanism checks \\
\bottomrule
\end{tabular}
\end{table}

\subsection{Model Configuration}

The Stage 1 LSTM-GPR model first learns temporal fuel-price patterns and then smooths the neural forecasts with Gaussian Process Regression. The GPR uses a constant kernel multiplied by a radial basis function kernel plus a white-noise term, with \texttt{alpha=0.01}, \texttt{n\_restarts\_optimizer=10}, and \texttt{random\_state=42}. Stage 2 uses an LSTM and feed-forward inflation predictor with an ensemble size of five. Stage 3 uses AutoGluon-Tabular with bagging and stacking enabled.

\begin{table}[!ht]
\centering
\caption{Stage 1 Price-Forecasting Settings}
\label{tab:hyperparameter-summary}
\label{tab:appendix-stage1-settings}
\scriptsize
\begin{tabular}{P{0.48\columnwidth}P{0.36\columnwidth}}
\toprule
\textbf{Parameter} & \textbf{Value} \\
\midrule
Time step & 12 months \\
Input scaling & Min-max, range (0, 1) \\
Conv1D filters / kernel & 80 / 3 \\
Bidirectional LSTM units & 128 \\
GRU units & 64 \\
Dense units & 80 \\
Dropout rates & 0.3, 0.5, 0.5 \\
Learning rate & 0.0005 \\
Epochs / batch size & 200 / 32 \\
Early stopping patience & 20 \\
GPR kernel & Constant $\times$ RBF + white noise \\
GPR optimizer restarts & 10 \\
GPR alpha / seed & 0.01 / 42 \\
\bottomrule
\end{tabular}
\end{table}

\begin{table}[!ht]
\centering
\caption{Stage 2 and Stage 3 Settings}
\label{tab:appendix-stage23-settings}
\scriptsize
\begin{tabular}{P{0.16\columnwidth}P{0.43\columnwidth}P{0.25\columnwidth}}
\toprule
\textbf{Stage} & \textbf{Parameter} & \textbf{Value} \\
\midrule
Stage 2 & Imputer neighbors & 3 \\
Stage 2 & Scaler & StandardScaler \\
Stage 2 & Lookback / time step & 4 \\
Stage 2 & Ensemble size & 5 \\
Stage 2 & LSTM units & 32 and 16 \\
Stage 2 & Dense units & 64 and 32 \\
Stage 2 & Regularization & L2(0.01) \\
Stage 2 & Epochs / batch size & 300 / 8 \\
Stage 2 & Validation split & 0.2 \\
Stage 3 & AutoGluon time limit & 600 seconds \\
Stage 3 & Bagging folds & 10 \\
Stage 3 & Stacking & auto stack \\
Stage 3 & GBM depth / learning rate & 3 / 0.01 \\
Stage 3 & CAT depth / learning rate & 3 / 0.01 \\
\bottomrule
\end{tabular}
\end{table}

\subsection{Forecasting and Robustness Outputs}

The archived oil-price comparison covers 47 observations in 2021--2024. Separate fixed-origin baseline forecasts use training data through 2020. The LSTM-GPR series has the smallest errors among the four archived comparison outputs. In Stage 3, the holdout split trains on 2026-01 to 2032-12 and tests on 2033-01 to 2035-12. Rolling-origin checks train through one calendar year and test the following year. These checks are used to separate predictive sensitivity from causal identification.

\begin{table}[!ht]
\centering
\caption{Stage 1 Oil-Price Forecast Metrics}
\label{tab:appendix-stage1-metrics}
\scriptsize
\begin{tabular}{lrrrr}
\toprule
\textbf{Model} & \textbf{$R^2$} & \textbf{MAPE} & \textbf{RMSE} & \textbf{MAE} \\
\midrule
LSTM-GPR & 0.963 & 0.0186 & 2.416 & 1.425 \\
LightGBM & 0.932 & 0.0299 & 3.267 & 2.464 \\
XGBoost & 0.729 & 0.0644 & 6.546 & 5.273 \\
SARIMA & -10.206 & 0.4839 & 42.069 & 39.449 \\
\bottomrule
\end{tabular}
\end{table}

\begin{table}[!ht]
\centering
\caption{Stage 3 Validation Summary}
\label{tab:appendix-validation}
\scriptsize
\begin{tabular}{lrrr}
\toprule
\textbf{Region} & \textbf{Holdout MAPE} & \textbf{Rolling MAPE} & \textbf{Positive Inc.} \\
\midrule
China & 0.0130 & 0.0100 & 8/27 \\
Europe & 0.0102 & 0.0078 & 6/27 \\
Middle East & 0.0555 & 0.0436 & 6/27 \\
\bottomrule
\end{tabular}
\vspace{2pt}

\raggedright\scriptsize
\textit{Note:} Holdout values use the Ridge inflation-trend or time-only specification with the lowest mean MAPE by region. Rolling values report the lowest mean MAPE across stable level/trend specifications. ``Positive Inc.'' counts models where adding fuel price and inflation improves out-of-sample $R^2$ over a time trend.
\end{table}

\subsection{Scenario Sensitivity and Transmission Diagnostics}

The price-shock table converts log-linear elasticities into average basis-point changes in WACC under a 10\% fuel-price shock. These values are conditional scenario sensitivities around the modeled WACC mean. The coefficient on log price is divided by the standard deviation of log price before conversion to original units. After reversing input standardization, the largest absolute responses in the cleaned panel occur in Middle East coal-price specifications. The table reports exact log-linear 10\% shock conversions around the modeled mean.

\begin{table}[!ht]
\centering
\caption{Largest 10\% Price-Shock Sensitivities}
\label{tab:appendix-elasticities}
\scriptsize
\begin{tabular}{llrr}
\toprule
\textbf{Region/Project} & \textbf{Fuel} & \textbf{WACC} & \textbf{bp} \\
\midrule
Middle East Onshore & Coal & 8.56\% & -26.058 \\
Middle East Solar & Coal & 7.54\% & -23.359 \\
Middle East Offshore & Coal & 7.60\% & -20.513 \\
Middle East Nuclear & Coal & 7.31\% & -14.954 \\
Middle East CCS & Coal & 7.54\% & -14.843 \\
Middle East Biomass & Coal & 7.66\% & -14.815 \\
Middle East Hydro & Coal & 6.52\% & -14.638 \\
China Offshore & Coal & 5.76\% & -7.500 \\
\bottomrule
\end{tabular}
\end{table}

\begin{table}[!ht]
\centering
\caption{Selected Granger-Style Transmission Diagnostics}
\label{tab:appendix-granger}
\scriptsize
\begin{tabular}{P{0.35\columnwidth}lrr}
\toprule
\textbf{Target} & \textbf{Fuel} & \textbf{Lag} & \textbf{$p$} \\
\midrule
China real interest rate & Coal & 1 & 0.00018 \\
China real interest rate & Coal & 2 & 0.00691 \\
Europe inflation & Gas & 2 & 0.00952 \\
China real interest rate & Gas & 1 & 0.01915 \\
China inflation & Coal & 1 & 0.09394 \\
\bottomrule
\end{tabular}
\vspace{2pt}

\raggedright\scriptsize
\textit{Note:} Diagnostics use first-differenced macro targets and log-differenced annual fuel prices. They indicate plausible transmission channels rather than structural causal identification.
\end{table}

\subsection{Replication Sequence}

To reproduce the manuscript results, first run the fuel-price preprocessing and prediction notebooks to generate price forecasts. Second, run the inflation prediction notebook to obtain macroeconomic covariates for the 2026-2035 horizon. Third, merge price, inflation, and WACC panels by region and project type and run the AutoGluon WACC evaluation. Finally, run the revision script to regenerate the holdout metrics, rolling-origin metrics, price-shock elasticities, and transmission diagnostics copied into the IEEE source tree. The raw WACC files are preserved, while cleaned regional inputs are written separately so the revision process remains auditable.
